
A modern világban a digitalizáció egyre inkább áthatja mindennapjainkat, a közlekedéstől kezdve a pénzügyi tranzakciókon át egészen az ügyintézésig. Az új technológiák bevezetése nem csupán kényelmesebbé, hanem hatékonyabbá és átláthatóbbá is teheti a különböző szolgáltatásokat. A tömegközlekedés területén azonban még mindig számos kihívással kell szembenéznünk, különösen a kisebb városokban és a regionális utazások során. Bár a nagyvárosokban egyre több helyen elérhető az elektronikus jegyrendszer és a valós idejű járatinformáció, sok vidéki térségben és távolsági közlekedési vonalon ezek a lehetőségek még mindig hiányoznak.

Az autóbuszok használata a közlekedés egyik leginkább környezetbarát formája, hiszen csökkenti a városi forgalmat és a szén-dioxid-kibocsátást. Ennek ellenére a buszközlekedés sokszor nem versenyképes alternatíva a magánautózással szemben, mivel a szolgáltatás minősége és elérhetősége gyakran nem felel meg a modern elvárásoknak. A jelenlegi rendszerek egyik legfőbb problémája a jegyvásárlási lehetőségek korlátozottsága: míg a városi közlekedésben már egyre inkább elérhető a bankkártyás és mobiltelefonos fizetés, a távolsági és regionális buszjáratok esetében sokszor még mindig csak készpénzzel lehet jegyet vásárolni a sofőrnél.

Egy másik fontos hiányosság, hogy sok esetben nem érhető el egy központi, megbízható információforrás a buszok menetrendjéről és aktuális helyzetéről. A buszmegállókban kifüggesztett menetrendek nem mindig naprakészek, és ha egy járat késik vagy kimarad, az utasoknak nincs lehetőségük arra, hogy erről időben értesüljenek. Egy digitális platform lehetőséget biztosíthatna a menetrendek valós idejű frissítésére, a buszok helyzetének követésére, valamint az utasok tájékoztatására késések esetén. Emellett sok felhasználó számára fontos lenne egy olyan alkalmazás, amely képes előre tervezni az utazásokat, például jegyfoglalási lehetőséggel vagy a járatok közötti átszállási lehetőségek figyelembevételével.

A jegyvásárlás és az utastájékoztatás digitalizálása nemcsak az utazók, hanem a busztársaságok számára is számos előnyt jelenthet. Az online jegyeladási rendszer segítségével pontosabb előrejelzések készíthetők az utasforgalomról, ami lehetővé teszi a járatok hatékonyabb üzemeltetését. Emellett a flottakezelési és adminisztrációs folyamatok is jelentősen egyszerűsíthetők egy digitális rendszer bevezetésével. Az autóbuszok műszaki állapotának és hivatalos dokumentumainak nyilvántartása papíralapon időigényes és kevésbé hatékony, míg egy online rendszer automatizált értesítésekkel segítheti a karbantartási és engedélyezési folyamatokat.

Fontos figyelembe venni azt is, hogy az utasok egyre nagyobb mértékben elvárják a digitális szolgáltatásokat, hiszen az okostelefonok és az internet folyamatos elérhetősége alapvetővé vált a mindennapi életben. Egy felhasználóbarát alkalmazás, amely egy helyen biztosítja az összes szükséges információt és lehetőséget a jegyvásárlásra, nemcsak kényelmesebbé teszi az utazást, hanem növeli is a tömegközlekedés iránti bizalmat. Az elmúlt években a digitális fizetési megoldások és az online utazástervezési eszközök iránti kereslet ugrásszerűen megnőtt, így a tömegközlekedési szolgáltatóknak is alkalmazkodniuk kell a változó utazási szokásokhoz.

A közlekedési rendszerek fejlesztése nem csupán technológiai kérdés, hanem társadalmi és környezetvédelmi szempontból is kiemelkedő jelentőséggel bír. A fenntartható közlekedési formák támogatása elengedhetetlen a nagyvárosi forgalmi dugók és a környezetszennyezés csökkentése érdekében. A buszközlekedés hatékonyabbá tétele és népszerűsítése hozzájárulhat a magánjármű-használat visszaszorításához, ezáltal mérsékelve a károsanyag-kibocsátást. Ha az emberek számára elérhetőbbé és megbízhatóbbá válik a tömegközlekedés, akkor nagyobb eséllyel választják azt a mindennapi utazásaik során, ami hosszú távon egy fenntarthatóbb városi mobilitást eredményezhet.

A rendszer két nagy részre osztható fel. Az egyik a frontend, amely a felhasználók számára látható felületeket, így a weboldalakat és a mobilalkalmazást tartalmazza, míg a másik rész a backend, amely az üzleti logikát, adatkezelést és API-végpontokat biztosítja. A backend részletes ismertetése az én dolgozatomban található meg, míg a frontend működését és implementációját külön dolgozat mutatja be \cite{bus4uui}.

Ebben a dolgozatban részletesen ismertetem a backend architektúrát, az API követelményeket, a technológiai választásokat, valamint az adatok tárolásának és kezelésének módját. Kitérünk a fejlesztés során alkalmazott tervezési döntésekre, a megvalósított funkciókra, továbbá a tesztelési módszerekre és a jövőbeni továbbfejlesztési lehetőségekre is. A dolgozat végén összegzem az eredményeket és levonom a projekt kapcsán szerzett tapasztalatokat.

